%% ma: L11-B26-0005
%% lop: 11
%% chuong: 7
%% bai: 26
%% dang: 11.26.3
%% loai: TLN
%% muc: VD
%% dap_an: "0,45"
%% nguon: "(NBV)-ĐỀ SỐ 10-MỖI NGÀY 1 ĐỀ THI 2025.pdf (D10, MỖI NGÀY 1 ĐỀ - Đề số 10), Phần III Câu 1"
%% kien_thuc: "Bài 26 (khoảng cách giữa hai đường thẳng chéo nhau), Bài 23 (đường thẳng vuông góc với mặt phẳng), định lí Pythagore"
%% trang_thai: chinh-thuc
%% ghi_chu: "Cùng ý tưởng với L11-B26-0003 (khoảng cách giữa hai đường chéo nhau trong hình chóp có đáy hình chữ nhật, cạnh bên hoặc đường cao đã biết), nhập vào dạng 11.26.3 đã duyệt. Gốc không có hình vẽ."
\begin{ex}
Cho hình chóp $S.ABCD$ có đáy $ABCD$ là hình vuông tâm $O$, cạnh $1$ cm. Cạnh bên $SA$ vuông góc với đáy, góc $\widehat{SBD}=60^\circ$. Khoảng cách giữa hai đường thẳng $AB$ và $SO$ bằng bao nhiêu cm? (Kết quả làm tròn đến hàng phần trăm.)
\shortans{0,45}
\loigiai{$ABCD$ là hình vuông cạnh $1$ nên $BD=\sqrt2$. Vì $SA\perp(ABCD)$ và $AB=AD$ nên $SB=SD$, do đó tam giác $SBD$ cân tại $S$ có $\widehat{SBD}=60^\circ$ nên đều, suy ra $SB=\sqrt2$ và $SA=\sqrt{SB^2-AB^2}=1$.

Gọi $M$, $N$ lần lượt là trung điểm của $AD$, $BC$; khi đó $MN\parallel AB$ và $O\in MN$. Do $AB\parallel(SMN)\ni SO$ nên $d(AB,SO)=d(AB,(SMN))=d(A,(SMN))$.

Vì $MN\parallel AB$ nên $MN\perp AD$, $MN\perp SA$, suy ra $MN\perp(SAD)$. Kẻ $AH\perp SM$ thì $AH\perp MN$, nên $AH\perp(SMN)$ và $d(A,(SMN))=AH$. Trong tam giác $SAM$ vuông tại $A$ với $SA=1$, $AM=\dfrac12$:
\[AH=\dfrac{SA\cdot AM}{\sqrt{SA^2+AM^2}}=\dfrac{\frac12}{\frac{\sqrt5}2}=\dfrac1{\sqrt5}\approx0{,}45\ \text{(cm)}.\]}
\end{ex}
